\section{Related Work}
%Logical Reasoning Dataset

%Hanmeng Liu et al. introduced LogiQA 2.0, an enhanced version of the original LogiQA dataset designed for evaluating logical reasoning~\cite{liu2023logiqa}. This updated version expands the dataset from 8,678 to 15,708 instances, receives high-quality English translations delivered by professional human translators, and includes a Natural Language Inference (NLI) dataset derived from the original MRC task. Experiments demonstrate that the increased dataset size and refined quality contribute to better performance across various pre-trained models, though state-of-the-art methods still fall short of human-level reasoning. The NLI section of LogiQA 2.0 is recognized as the first large-scale, expert-designed NLI benchmark for logical reasoning.
%
%ZX: If we do not use this dataset, LogiQA dataset, why did we mention it here? This gives the reviewers an opportunity to reject our paper for including irrelevant research. 
%

%Situation Puzzles Evaluation

%Tian Zhao et al. proposed SPQA, a new question-answering model specifically designed to tackle situation puzzle questions by incorporating a novel situation-puzzle related dataset (SpQ)~\cite{zhao2023question}. Built on the T5 architecture similar to UnifiedQA, SPQA is fine-tuned and prompt-tuned to assess its performance in handling both situation puzzle questions and standard yes/no questions. Their study reveal fine-tuning with the SpQ dataset significantly enhances the model's ability to solve situation puzzles but reduces its performance on general yes/no questions. Prompt-tuning further highlights the importance of the SpQ dataset, with a notable improvement in solving situation puzzles.
%
%%ZX: When the reviewers read the paragraph above, they will expect us to use the situation-puzzle related dataset for comparison. However, our paper did not use the dataset. Why don’t we consider the situation-puzzle related dataset you mentioned above? It would be interesting to use the same dataset that other researchers have utilized.
%
%OpenAI o1
%
%OpenAI o1 is a cutting-edge model that excels in competitive programming and complex reasoning tasks. It ranks in the 89th percentile on Codeforces and achieves top-tier performance on the USA Math Olympiad qualifier (AIME), solving up to 93\% of problems using consensus-based sampling techniques. OpenAI o1 also surpasses human PhD-level performance on the GPQA-diamond benchmark, which assesses expertise in biology, physics, and chemistry. Notably, it is the first model to be competitive with human experts on MMLU's MMMU subset, scoring 78.2\%. Leveraging a chain-of-thought strategy, o1 utilizes reinforcement learning to break down complex problems into simpler steps, continuously refining its reasoning and problem-solving capabilities. This advanced training method allows the model to improve both its accuracy and efficiency over time.
%%
% ZX: Detailed information from OpenAI's report above provides more context. The statistic stating that '93% of problems use consensus-based sampling techniques' needs to be properly cited; please specify the source. Since I could not find the exact numbers, I refrained from using them. Please feel free to add it into Related Work Section (e.g., second pragraph).
%
%
%Qwen2.5
%The Qwen2.5 series models represent the newest frontier in the Qwen family. They are trained on 18 trillion tokens and surpasses Qwen2 significantly in knowledge (MMLU: 85+), as well as proficiency in coding (HumanEval 85+) and mathematics (MATH 80+). The flagship model Qwen2.5-72B outperforms the newly released Llama3 70B model and the mixture of experts Mixtral 8x22B model across various benchmarks, and remains highly competitive even when compared against larger models like Llama-3-405B\cite{qwen2.5}.  On the other end of the spectrum, Qwen2.5-3B demonstrates the accelerated growth in knowledge density among language models is a prime example of small language models (SLMs) reducing the gap with LLMs. Additionally, the new models achieve significant improvements in instruction following, generating long texts (over 8K tokens), understanding structured data (e.g, tables), and generating structured outputs especially JSON. Qwen2.5 models are generally more resilient to the diversity of system prompts, enhancing role-play implementation and condition-setting for chatbots.   
%zx: rewrite as below in the related work section.

%Ensemble Model 
%
%Chengyan Wu et al. (2024) proposed an ensemble model for diabetes question classification, combining multiple large and small language models to enhance performance. The model ensemble includes Qwen-7b-Chat, ChatGLM2-6b, and MacBERT as base models, fine-tuned using the QLoRA and LoRA frameworks for instruction-based learning. These large models focus on multiple-choice tasks, while smaller models trained with adversarial techniques, such as Fast Gradient Method (FGM), handle traditional text classification. To further improve accuracy, misclassified samples are identified and augmented via paraphrasing, adding diversity to the training data. A majority voting mechanism is then applied to aggregate predictions from the models, ensuring a robust final classification. Their findings suggest that integrating advanced
%language models with few-shot learning strategies can significantly enhance classification performance, especially when the ground truth data is generated or verified by the same type of models.
% ZX: why we introduce Ensemble Model here? it is irrelevant to this paper.
%
%
The rapid evolution of large language models (LLMs) has generated significant interest in their applications for complex reasoning tasks \cite{chen2024large, hadi2024large}. Early research underscored the foundational capabilities of models like GPT-3 and its successors in handling basic logical reasoning and natural language processing tasks \cite{zhang2021commentary, kalyan2023survey}. While these models excelled in pattern recognition and could infer information from extensive datasets, challenges remained, particularly in maintaining coherent reasoning across multiple steps \cite{yang2024harnessing}.

Recent advancements, such as OpenAI's GPT-4, have enhanced deductive reasoning abilities through techniques like chain-of-thought prompting, which encourage models to articulate their reasoning processes more clearly \cite{gou2024rationality}. Despite these improvements, LLMs still exhibit limitations in multi-step reasoning tasks, often leading to cumulative errors in lengthy logical chains \cite{tong2024can}.

The landscape of LLMs is rapidly expanding, with significant contributions from both proprietary and open-source models \cite{he2023survey}. Meta's Llama series and newer models from Alibaba, such as Qwen2.5, have emerged as strong competitors to established leaders \cite{hui2024qwen2, yang2024qwen2}. These models not only offer large context windows and multilingual support but also incorporate innovative features like Grouped Query Attention, enhancing efficiency in processing complex queries \cite{yang2024harnessing}.

% The Qwen2.5 series represents the latest advancements in the Qwen family. Trained on 18 trillion tokens, these models significantly surpass Qwen2 in knowledge (MMLU: 85+), coding proficiency (HumanEval: 85+), and mathematical abilities (MATH: 80+). The flagship model, Qwen2.5-72B, outperforms the recently released Llama3 70B model and the mixture of experts, Mixtral 8x22B, across various benchmarks, remaining highly competitive even against larger models like Llama-3-405B \cite{qwen2.5}. Conversely, Qwen2.5-3B exemplifies the accelerated growth in knowledge density among language models, demonstrating how small language models (SLMs) are narrowing the performance gap with larger language models (LLMs). Furthermore, these new models achieve substantial improvements in instruction following, generating long texts (over 8K tokens), understanding structured data (e.g., tables), and producing structured outputs, particularly in JSON format. The Qwen2.5 models are generally more resilient to diverse system prompts, enhancing role-play implementation and condition-setting for chatbots.

Moreover, the performance of Chinese LLMs, such as InternLM2.5, reflects a growing emphasis on addressing cultural and linguistic nuances within AI systems. These models are specifically designed to accommodate bilingual contexts and align with the unique social and historical aspects of Asian cultures, providing potential advantages in cross-cultural applications \cite{zhang2024mme}.

Several studies have begun to investigate the competitive landscape of LLMs in reasoning tasks. Recent evaluations have focused on few-shot and fine-tuning methodologies to determine optimal performance metrics across various architectures \cite{winata2021language}. However, the exploration of logical reasoning in contexts like the Turtle Soup game remains limited, presenting opportunities for further investigation.

% This study builds on the existing body of research by systematically comparing the latest iterations of LLMs against proprietary models in complex reasoning scenarios, utilizing both few-shot prompting and fine-tuning techniques. By introducing the Valid Classification Ratio (VCR) as a new metric for assessing model output reliability, our work aims to contribute valuable insights into the evolving capabilities of LLMs in logical reasoning tasks.
